\documentclass[11pt,a4paper]{ctexart}
\usepackage[margin=2.2cm]{geometry}
\usepackage{xcolor}
\usepackage{enumitem}
\usepackage{amssymb}
\usepackage{booktabs}
\usepackage[normalem]{ulem}
\usepackage{array}
\usepackage{longtable}
\usepackage{titlesec}
\usepackage{hyperref}
\definecolor{navy}{HTML}{1A365D}
\hypersetup{colorlinks=true,linkcolor=navy,urlcolor=navy}
\setlist{nosep}
\titleformat{\section}{\Large\bfseries\color{navy}}{}{0pt}{}
\titleformat{\subsection}{\large\bfseries\color{navy}}{}{0pt}{}
\pagestyle{plain}
\renewcommand{\arraystretch}{1.25}

\title{\textbf{七年级上册语法总结}\\[4pt]
\large 人教版（新课纲版）七年级上册 · 按教材书末 Grammar 附录整理}
\author{初中英语资料站}
\date{2026-11-12 · 附配套专项练习 192 题}
\begin{document}
\maketitle
\section*{使用说明}
\begin{itemize}[leftmargin=1.4em]
\item 本讲义按七年级上册教材书末 Grammar 附录整理，表格与例句取自教材原文。
\item 网页版：\url{https://1901151458-prog.github.io/english-materials/g7a-grammar/}
\end{itemize}

\section{词类（Parts of Speech）}
\noindent\textcolor{gray}{英语的九大词类}\par\smallskip
{\footnotesize
\begin{longtable}{p{3.35cm}p{3.35cm}p{3.35cm}p{3.35cm}}
\hline
词类 & 英语名称 & 意义 & 例词 \\
\hline
\endfirsthead
\hline
词类 & 英语名称 & 意义 & 例词 \\
\hline
\endhead
\hline
\endfoot
名词 & Noun (n.) & 表示人或事物的名称 & son, tennis, family, idea \\
冠词 & Article (art.) & 用于名词前，帮助限定名词所指的人或事物 & a, an, the \\
代词 & Pronoun (pron.) & 用于代替名词以及起名词作用的短语、分句或句子等 & we, her, that, what \\
形容词 & Adjective (adj.) & 用于修饰名词，表示人或事物的特征 & funny, different, excited, amazing \\
数词 & Numeral (num.) & 表示数量或顺序 & nine, thirteen, first, twentieth \\
动词 & Verb (v.) & 表示动作或状态 & be (am, is, are), play, think, have \\
副词 & Adverb (adv.) & 用于修饰动词、形容词或其他副词 & not, too, there, usually \\
介词 & Preposition (prep.) & 表示名词、代词等与句中其他词的关系 & in, on, at, of \\
连词 & Conjunction (conj.) & 用于连接单词、短语或句子 & and, but, because, when \\
感叹词 & Interjection (interj.) & 表示说话时的喜悦、惊讶等情感 & well, hello, oh, hey \\
\end{longtable}
}

\begin{center}\fbox{\parbox{0.92\textwidth}{\small \textbf{说明：}教材表1 实际列出 10 个名称（含感叹词），其中「冠词」是虚词、其余多为实词。记住每个词类的\textbf{英文缩写}（n. / v. / adj. / adv. / prep. / conj. / pron. / num. / art. / interj.）是查词典的基础。}}\end{center}

\section{名词（Nouns）}
\noindent\textcolor{gray}{种类、可数性、所有格}\par\smallskip
名词是指人或事物的名称。

\subsection*{1. 名词的种类}

名词可分为\textbf{专有名词}和\textbf{普通名词}两类。专有名词是特定的人、事件、机构、地点等专有的名称，如 Peter、January、Singapore。\textbf{专有名词的第一个字母一般要大写}。

\subsection*{2. 名词的可数性}

名词按其可数性可分为\textbf{可数名词}和\textbf{不可数名词}。可数名词有复数形式，如 eight rabbits、some trees。不可数名词一般没有复数形式，如 milk、bread、yogurt。

\subsection*{3. 名词所有格}

名词所有格表示所属关系。

{\footnotesize
\begin{longtable}{p{3.35cm}p{3.35cm}p{3.35cm}p{3.35cm}}
\hline
类别 & 构成方法 & 例词 & 读音 \\
\hline
\endfirsthead
\hline
类别 & 构成方法 & 例词 & 读音 \\
\hline
\endhead
\hline
\endfoot
单数名词 & 加 \textbf{’s} & Mike’s basketball / your mum’s birthday / Kate’s cat / Alice’s vlog / the actress’s name / George’s father & 在清辅音后读 /s/，在浊辅音和元音后读 /z/，在 /s/、/z/、/ʃ/、/tʃ/、/dʒ/ 等后读 /ɪz/ \\
复数名词（不以 -s 结尾） & 加 \textbf{’s} & Children’s Day & 同上 \\
复数名词（以 -s 结尾） & 加 \textbf{’} & her grandparents’ favourite grandchild & 读音不变 \\
\end{longtable}
}

\begin{center}\fbox{\parbox{0.92\textwidth}{\small \textbf{口诀：}单数加 ’s；复数词尾有 s 只加 ’；复数词尾无 s 仍加 ’s。}}\end{center}

\section{冠词（Articles）}
\noindent\textcolor{gray}{定冠词 the、不定冠词 a/an、零冠词}\par\smallskip
冠词是置于名词之前、限定名词所表示的人或事物的一种虚词。冠词不能离开名词单独使用。分为\textbf{定冠词（the）}、\textbf{不定冠词（a / an）}和\textbf{零冠词}。

\subsection*{1. 定冠词 the}

定冠词的基本概念是「\textbf{特指}」，用于指说话人和听话人已知的人或事物，可与单数可数名词、复数可数名词或不可数名词一起使用。

\begin{itemize}[leftmargin=1.4em]
\item Where is \underline{the} library?\\ 图书馆在哪儿？（双方都知道的那个）
\item \underline{The} glasses are on her dad’s nose.\\ 眼镜在她爸爸的鼻子上。
\item \underline{The} milk is on the table.\\ 牛奶在桌上。
\end{itemize}

定冠词在辅音前读 /ðə/，在元音前读 /ði/，特别强调时可读作 /ðiː/。

\subsection*{2. 不定冠词 a / an}

不定冠词的基本概念是「\textbf{非特指}」，用于指某类人或物中的任何一个或某一个。\textbf{an 放在以元音（音素）开头的名词前面}，如 an eraser。不定冠词只能与单数可数名词一起使用，其基本含义有「该类中的一例」「只有一个」「每一个」等。

\begin{itemize}[leftmargin=1.4em]
\item I have \underline{a} pet dog.\\ 我有一只宠物狗。（一只）
\item Sometimes I play basketball for \underline{an} hour.\\ 我有时打一小时篮球。（一个小时）
\item We have three lessons \underline{a} day.\\ 我们一天上三节课。（每一）
\end{itemize}

\subsection*{3. 零冠词}

零冠词指名词前不加冠词的情况。

\begin{itemize}[leftmargin=1.4em]
\item My favourite subject is \underline{English}.\\ 我最喜欢的科目是英语。
\item Binbin is good with \underline{numbers}.\\ 彬彬擅长数字。
\item It’s time for me to go to \underline{bed}.\\ 我该上床睡觉了。
\end{itemize}

\begin{center}\fbox{\parbox{0.92\textwidth}{\small \textbf{教材注释：}有的语法书并不把零冠词列为冠词中的一种。}}\end{center}

\section{代词（Pronouns）}
\noindent\textcolor{gray}{人称代词、形容词性物主代词、指示代词}\par\smallskip
代词是代替名词以及起名词作用的短语、分句或句子的词，它们的词义必须通过上下文来确定。

\subsection*{1. 人称代词}

{\footnotesize
\begin{longtable}{p{2.68cm}p{2.68cm}p{2.68cm}p{2.68cm}p{2.68cm}}
\hline
人称 & 单数主格 & 单数宾格 & 复数主格 & 复数宾格 \\
\hline
\endfirsthead
\hline
人称 & 单数主格 & 单数宾格 & 复数主格 & 复数宾格 \\
\hline
\endhead
\hline
\endfoot
第一人称 & I（我） & me & we（我们） & us \\
第二人称 & you（你） & you & you（你们） & you \\
第三人称 & he（他）\\ she（她）\\ it（它） & him\\ her\\ it & they（他们/她们/它们） & them \\
\end{longtable}
}

\subsection*{2. 形容词性物主代词}

物主代词是表示所有关系的代词，分为形容词性物主代词和名词性物主代词。此处重点介绍\textbf{形容词性物主代词}，它们相当于形容词，置于名词之前作定语，如 my name、your schoolbag、her favourite subject、their children。

{\footnotesize
\begin{longtable}{p{4.47cm}p{4.47cm}p{4.47cm}}
\hline
人称 & 单数 & 复数 \\
\hline
\endfirsthead
\hline
人称 & 单数 & 复数 \\
\hline
\endhead
\hline
\endfoot
第一人称 & my & our \\
第二人称 & your & your \\
第三人称 & his / her / its & their \\
\end{longtable}
}

\subsection*{3. 指示代词}

{\footnotesize
\begin{longtable}{p{4.47cm}p{4.47cm}p{4.47cm}}
\hline
指示代词 & 用法 & 例句 \\
\hline
\endfirsthead
\hline
指示代词 & 用法 & 例句 \\
\hline
\endhead
\hline
\endfoot
this（这个）\\ these（这些） & 用于指时间或空间上\textbf{较近}的事物 & This is my uncle’s farm.\\ These are my friends at the drama club. \\
that（那个）\\ those（那些） & 用于指时间或空间上\textbf{较远}的事物 & Do you see that red house by the river?\\ Those ping-pong bats in the corner are my grandpa’s. \\
\end{longtable}
}

\section{数词（Numerals）}
\noindent\textcolor{gray}{基数词与序数词}\par\smallskip
\subsection*{1. 基数词}

基数词表示数目的多少。

{\footnotesize
\begin{longtable}{p{3.35cm}p{3.35cm}p{3.35cm}p{3.35cm}}
\hline
1–10 & 11–19 & 整十 & 其他 \\
\hline
\endfirsthead
\hline
1–10 & 11–19 & 整十 & 其他 \\
\hline
\endhead
\hline
\endfoot
one, two, three, four, five, six, seven, eight, nine, ten & eleven, twelve, thirteen, fourteen, fifteen, sixteen, seventeen, eighteen, nineteen & twenty, thirty, forty, fifty, sixty, seventy, eighty, ninety & twenty-one, twenty-two, one hundred, one hundred and one, two hundred \\
\end{longtable}
}

\subsection*{2. 序数词}

序数词表示顺序的先后。\textbf{序数词前常用定冠词 the 或物主代词}，有时也可用不定冠词或零冠词。

{\footnotesize
\begin{longtable}{p{3.35cm}p{3.35cm}p{3.35cm}p{3.35cm}}
\hline
1st–10th & 11th–19th & 整十 & 特殊 \\
\hline
\endfirsthead
\hline
1st–10th & 11th–19th & 整十 & 特殊 \\
\hline
\endhead
\hline
\endfoot
first, second, third, fourth, fifth, sixth, seventh, eighth, ninth, tenth & eleventh, twelfth, thirteenth, fourteenth, fifteenth, sixteenth, seventeenth, eighteenth, nineteenth & twentieth, thirtieth, fortieth, fiftieth, sixtieth, seventieth, eightieth, ninetieth & twenty-first (21st), twenty-second (22nd), one hundredth (100th), one hundred and first (101st) \\
\end{longtable}
}

\begin{center}\fbox{\parbox{0.92\textwidth}{\small \textbf{记忆要点：}① first / second / third 不规则；② fifth、eighth、ninth、twelfth 拼写特殊；③ 整十把 y 变 ie 加 th；④ 几十几的序数词只把\textbf{个位}变序数（twenty-first）。}}\end{center}

\section{动词（Verbs）}
\noindent\textcolor{gray}{系动词、实义动词、助动词、情态动词}\par\smallskip
动词表示动作或状态，人称和数一般必须与主语一致。此处简要介绍\textbf{系动词、实义动词、助动词和情态动词}。

\subsection*{1. 系动词}

系动词是连接主语和表语的动词，不能独立作谓语，必须与其后的表语一起构成谓语。常见的系动词有 \textbf{be（am、is、are）、look、get、turn} 等。

\begin{itemize}[leftmargin=1.4em]
\item I \underline{am} 12 years old.\\ 我 12 岁。
\item You \underline{look} happy today.\\ 你今天看起来很开心。
\item Trees \underline{turn} green in spring.\\ 春天树木变绿。
\end{itemize}

\subsection*{2. 实义动词}

实义动词意义完全，能独立作谓语。

\begin{itemize}[leftmargin=1.4em]
\item I \underline{play} tennis with my friends after school.\\ 放学后我和朋友打网球。
\item My father \underline{likes} fishing a lot.\\ 我爸爸很喜欢钓鱼。
\end{itemize}

\subsection*{3. 助动词}

助动词本身无词汇意义或意义不完全，不能独立作谓语，只能和实义动词一起构成各种时态、语态、语气，以及否定和疑问结构。常见的助动词有 \textbf{be、do、have、will} 等。

\begin{itemize}[leftmargin=1.4em]
\item \underline{Do} you want to be my friend?\\ 你想成为我的朋友吗？
\item I \underline{don’t} have a brother.\\ 我没有兄弟。
\end{itemize}

\subsection*{4. 情态动词}

情态动词表示说话人的语气和情态，不能独立作谓语，只能和主要动词一起构成谓语。常见的有 \textbf{can、may、could、must、have to、should、would} 等，\textbf{其后动词使用原形}。情态动词一般没有人称和数的变化，否定式一般在其后加 not，一般疑问句中通常提到句首。

本册主要学习情态动词 \textbf{can 表示「能力」}的用法：

{\footnotesize
\begin{longtable}{p{6.70cm}p{6.70cm}}
\hline
句式 & 例句 \\
\hline
\endfirsthead
\hline
句式 & 例句 \\
\hline
\endhead
\hline
\endfoot
肯定式 & Linda \textbf{can swim} really well.\\ Cats \textbf{can climb} trees. \\
否定式 & I \textbf{can’t (cannot) cook}.\\ My parents \textbf{can’t (cannot) speak} English. \\
疑问式和回答 & \textbf{Can you play} ping-pong? — Yes, I can. / No, I can’t.\\ \textbf{What instruments can} Teng Fei \textbf{play}? — He can play the drums. \\
\end{longtable}
}

\section{介词（Prepositions）}
\noindent\textcolor{gray}{时间介词、地点介词、固定搭配}\par\smallskip
介词一般用于名词或代词前，表示该词与句中其他成分的关系。介词后面的名词或代词称为介词宾语。\textbf{介词宾语若是人称代词，则要用宾格}。介词和介词宾语合在一起构成介词短语。

\subsection*{1. 表示时间的介词 at / on / in}

{\footnotesize
\begin{longtable}{p{4.47cm}p{4.47cm}p{4.47cm}}
\hline
at & on & in \\
\hline
\endfirsthead
\hline
at & on & in \\
\hline
\endhead
\hline
\endfoot
钟点 at 3 o’clock\\ 用餐 at lunchtime\\ 节日 at Christmas\\ 年龄 at the age of 20\\ 时间 at the same time\\ 一天中的某段时间 at night & 日期 on 22nd November\\ 星期几 on Thursday\\ 节假日 on New Year’s Day\\ 某一天的某段时间 on Sunday mornings & 月份 in April\\ 季节 in spring\\ 年份 in 1949\\ 一天中的某段时间 in the afternoon \\
\end{longtable}
}

其他表示时间的介词还有 \textbf{before、after} 等，如 before breakfast、after school。

\subsection*{2. 表示地点的介词}

\paragraph{at / on / in}

\begin{itemize}[leftmargin=1.4em]
\item \textbf{at} 表示在某个地点：at school、at home、at the bus stop
\item \textbf{on} 表示在某个表面上：on the desk、on the door、on the sports field
\item \textbf{in} 表示在某个范围内：in the photo、in the garden、in China、in Class 2
\end{itemize}

\paragraph{方位介词}

\begin{itemize}[leftmargin=1.4em]
\item \textbf{under} 在……下面
\item \textbf{in front of} 在……前面
\item \textbf{behind} 在……后面
\item \textbf{next to} 紧邻；在……近旁
\item \textbf{across from} 在……对面
\item \textbf{between} 在……中间
\end{itemize}

\subsection*{3. 与其他词的连用}

{\footnotesize
\begin{longtable}{p{6.70cm}p{6.70cm}}
\hline
搭配类型 & 例子 \\
\hline
\endfirsthead
\hline
搭配类型 & 例子 \\
\hline
\endhead
\hline
\endfoot
介词与名词 & on duty 值日；值班 / on the way to 在去……的路上 / on the farm 在农场 / on the left 在左面 / in the middle 在中间 / in the future 将来 \\
动词与介词 & look at 看 / talk about 谈论 / think of 想到 / learn from 向……学习 / listen to 听 / ask for 请求 \\
形容词与介词 & be good at 在某方面擅长 / be famous for 因……而出名 / be similar to 与……相似 / be interested in 对……感兴趣 \\
\end{longtable}
}

\section{一般现在时（Simple Present Tense）}
\noindent\textcolor{gray}{意义、构成、三单变化}\par\smallskip
\subsection*{1. 一般现在时的意义}

\begin{itemize}[leftmargin=1.4em]
\item 表示\textbf{现在的状态}：Where are the twins from? / She doesn’t have a brother.
\item 表示\textbf{经常发生的或习惯性的动作}：My school begins at 9:00. / I often listen to music on the way.
\item 表示\textbf{主语具备的性格和能力}等：Colin’s father works hard. / I like all my classes, but my favourite is maths.
\end{itemize}

\subsection*{2. 一般现在时的构成}

\paragraph{（1）系动词 be}

{\footnotesize
\begin{longtable}{p{6.70cm}p{6.70cm}}
\hline
肯定式 & 否定式 \\
\hline
\endfirsthead
\hline
肯定式 & 否定式 \\
\hline
\endhead
\hline
\endfoot
I am … / You are … / He / She / It is … / We / You / They are … & I am not … / You are not … / He / She / It is not … / We / You / They are not … \\
Am I …? / Are you …? / Is he / she / it …? / Are they …? & Yes, you are. / Yes, I am. / Yes, he / she / it is. / Yes, they are. \\
\end{longtable}
}

{\footnotesize
\begin{longtable}{p{4.47cm}p{4.47cm}p{4.47cm}}
\hline
肯定式 & 否定式 & 其他 \\
\hline
\endfirsthead
\hline
肯定式 & 否定式 & 其他 \\
\hline
\endhead
\hline
\endfoot
I’m = I am / you’re = you are / he’s = he is / she’s = she is / it’s = it is / we’re = we are / they’re = they are & isn’t = is not / aren’t = are not & that’s = that is / what’s = what is / who’s = who is / who’re = who are / where’s = where is / where’re = where are \\
\end{longtable}
}

\paragraph{（2）实义动词}

{\footnotesize
\begin{longtable}{p{6.70cm}p{6.70cm}}
\hline
肯定式 & 否定式 \\
\hline
\endfirsthead
\hline
肯定式 & 否定式 \\
\hline
\endhead
\hline
\endfoot
I / You like oranges.\\ He / She \textbf{likes} oranges.\\ We / You / They like oranges. & I / You \textbf{do not} like oranges.\\ He / She \textbf{does not} like oranges.\\ We / You / They do not like oranges. \\
\textbf{Do} you like oranges?\\ \textbf{Does} he / she like oranges? & Yes, I do. / No, I do not.\\ Yes, he / she does. / No, he / she does not. \\
\end{longtable}
}

\textbf{注：}在口语中，do not 经常缩略为 \textbf{don’t}，does not 经常缩略为 \textbf{doesn’t}。

{\footnotesize
\begin{longtable}{p{4.47cm}p{4.47cm}p{4.47cm}}
\hline
类别 & 构成方法 & 例词 \\
\hline
\endfirsthead
\hline
类别 & 构成方法 & 例词 \\
\hline
\endhead
\hline
\endfoot
一般情况 & 加 \textbf{-s} & help → helps / like → likes / come → comes / know → knows / play → plays / get → gets / find → finds \\
以 -s、-x、-ch、-sh 结尾 & 加 \textbf{-es} & guess → guesses / fix → fixes / teach → teaches / finish → finishes \\
以辅音字母加 -y 结尾 & 变 y 为 i，再加 \textbf{-es} & study → studies \\
\end{longtable}
}

\begin{center}\fbox{\parbox{0.92\textwidth}{\small \textbf{读音规则：}在清辅音后读 /s/，在浊辅音和元音后读 /z/，在 /s/、/z/、/ʃ/、/tʃ/、/dʒ/ 等后读 /ɪz/。}}\end{center}

\section{句子种类（Sentence Types）}
\noindent\textcolor{gray}{陈述句与疑问句}\par\smallskip
英语句子按照用途可分为\textbf{陈述句、疑问句、祈使句和感叹句}四类。本册重点介绍\textbf{陈述句和疑问句}。

\subsection*{1. 陈述句}

陈述句包括肯定结构和否定结构。在肯定结构中，谓语动词不含否定词；在否定结构中，系动词 be、助动词或情态动词后加 not（常用缩略形式），或用 no、never 等词表示。陈述句末用句号。

{\footnotesize
\begin{longtable}{p{6.70cm}p{6.70cm}}
\hline
肯定式 & 否定式 \\
\hline
\endfirsthead
\hline
肯定式 & 否定式 \\
\hline
\endhead
\hline
\endfoot
Ella’s cap \textbf{is} in her schoolbag. & Emma’s cap \textbf{is not} in her schoolbag. \\
He \textbf{has} some ping-pong bats. & He \textbf{doesn’t have} any ping-pong balls. \\
They \textbf{usually take} showers in the morning. & They \textbf{never take} showers in the afternoon. \\
\end{longtable}
}

\begin{center}\fbox{\parbox{0.92\textwidth}{\small \textbf{否定三种方式：}① be 动词后加 not（is not）；② 实义动词用 don’t / doesn’t + 动词原形；③ 用 never / no 等否定词。}}\end{center}

\section{配套专项练习}
\noindent\textcolor{gray}{四个专项各 48 题、共 192 题，支持在线作答与即时解析}\par\smallskip
\paragraph{名词与所有格专项 · 48 题}

\paragraph{冠词与代词专项 · 48 题}

\paragraph{数词与介词专项 · 48 题}

\paragraph{一般现在时与句子种类专项 · 48 题}

\begin{center}\fbox{\parbox{0.92\textwidth}{\small \textbf{建议顺序：}先按本页 01–09 讲通规则 → 再进入对应专项练习自测 → 错题回到小节复习。四个专项合计 \textbf{192 题}。\\ \textbf{考前速查：}一页纸速查表 →（A4 单页，可打印）}}\end{center}

\end{document}
